\chapter{Introduction}

Machine learning for wearable Human Activity Recognition (HAR) is well studied in controlled experimental settings, but many systems still fail to translate reliably into production-like operation. In practice, data distributions change, labels are unavailable at inference time, and model quality can degrade silently when sensor behavior, user context, or preprocessing assumptions shift. This gap between offline model performance and operational robustness creates a core engineering challenge: achieving trustworthy model behavior over time, not only high benchmark accuracy.

This thesis addresses that challenge through an offline-first MLOps perspective for HAR. The goal is not only to train an accurate classifier, but to design an end-to-end pipeline that remains reproducible, monitorable, and maintainable under realistic constraints. In particular, the system must support repeatable data and model workflows, continuous post-inference quality checks, drift-aware monitoring, and trigger-based retraining decisions. Prior software engineering research for machine learning systems emphasizes that production reliability depends on disciplined lifecycle practices beyond model code alone \cite{amershi2019mlops}. That principle is central to this work.

The problem setting for this thesis is intentionally practical. The repository integrates data ingestion, validation, preprocessing, model training, inference, evaluation, monitoring, and retraining governance into one orchestrated flow. This enables evidence-driven assessment of whether the system is operating correctly and whether observed signals should lead to no action, warning, or retraining actions. The contribution is therefore both technical and methodological: a concrete implementation and an auditable workflow that links model behavior to operational decisions.

The rest of this thesis is organized as follows. Chapter 2 summarizes the required background and related work. Chapter 3 presents the offline MLOps architecture. Chapter 4 describes data processing and model development. Chapter 5 details implementation choices in the repository. Chapter 6 focuses on monitoring, drift detection, and retraining logic. Chapter 7 reports system experiments. Chapter 8 concludes with limitations and future work.
